\section{Diffusion-Based Style Transfer}
\label{sec:style-transfer}

This section addresses \textbf{RQ4}: it targets the domain and Sim2Real gap by
transferring concrete style components of the target domain onto the labeled synthetic
source --- the \texttt{combined} dataset --- thereby injecting domain information. It
is the most direct appearance-level intervention: rather than approximating weather
with library effects (\Cref{sec:weather-aug}) or relying on the model's own
predictions (\Cref{sec:pseudo-labels}), it repaints the source images to look like the
target. This stage is \emph{planned/ongoing}.

\subsection{Motivation and Label Preservation}
\label{subsec:st-motivation}
Style transfer imports stylistic properties of the real domain --- textures, colour
contrasts --- and weather/time-of-day scenarios that either look different or are
absent in the synthetic data. The style gap quantified in \Cref{subsec:gap-weather}
(the distinct ``real'' cluster in the UMAP of style features) is precisely what this
stage aims to close. The core constraint is \emph{label preservation}: the four runway
corners must remain geometrically valid after restyling. CycleGAN-style unpaired
translation can change texture and colour but is hard to control, and if the
transformation displaces the four corners the synthetic labels become incorrect.

\subsection{Method: Control-Guided Diffusion (Cosmos Transfer)}
\label{subsec:st-method}
Recent diffusion approaches use a \emph{control branch} that steers the diffusion
process. This branch is an additional model that combines the source image with
\emph{control modalities} such as text prompts, reference images or
segmentation/edge maps. By conditioning generation on the source structure, the
original labels are preserved far better than with GAN-based approaches. The thesis
uses the open-source \textbf{Cosmos Transfer} diffusion model, conditioned on the
structure of the \texttt{combined} synthetic images so that real and adverse-weather
appearance is applied while the runway geometry --- and hence the 4-corner labels ---
is retained. The restyled images are then used to fine-tune YOLOv8-pose under the full
budget.
\todo[inline]{Fix the exact Cosmos Transfer version and the control modality used
(edge/segmentation/depth), the target styles/weather applied, prompt settings, and the
number of restyled images.}

\begin{table}[th]
    \caption{Planned evaluation: control-guided diffusion style transfer applied to
    \texttt{combined}, on the real test set per weather scenario, relative to the
    \texttt{combined} baseline.}
    \label{tab:st-results}
    \centering
    \begin{NiceTabular}{lrrrrrr}
        \CodeBefore
        \rowcolors{2}{gray!10}{white}
        \Body
        \toprule
        \textbf{Configuration} & \textbf{Nominal} & \textbf{Glare} & \textbf{Rain} & \textbf{Fog} & \textbf{Night} & \textbf{Overall} \\
        \midrule
        combined (baseline)  & \tbd & \tbd & \tbd & \tbd & \tbd & \tbd \\
        + Cosmos Transfer    & \tbd & \tbd & \tbd & \tbd & \tbd & \tbd \\
        \bottomrule
    \end{NiceTabular}
\end{table}

\plotslot{Style-transfer examples: synthetic source, edge/segmentation control map,
and diffusion-restyled output with preserved runway geometry.}

\subsection{Optional Comparison: CycleGAN Baseline}
\label{subsec:st-cyclegan}
Time permitting, the diffusion results are compared against a CycleGAN baseline that
applies the same source$\rightarrow$target translation but may displace the corners.
The comparison quantifies both the downstream accuracy and the label-integrity
advantage of the control-guided diffusion approach (the corner displacement introduced
by the translation itself).

\subsection{Discussion}
\label{subsec:st-discussion}
Control-guided diffusion aims to combine the strengths of the earlier techniques: like
augmentation it keeps labels exact and needs no target labels, and like self-training
it grounds the adaptation in real target appearance --- while, unlike CycleGAN, it is
steerable and structure-preserving. Its main costs are compute and the fidelity and
controllability of the generative model. The cross-technique comparison in
\Cref{sec:discussion} will place these results alongside weather augmentation and
self-training once available.
